\chapter{Estado del arte}
\label{cap:sota}

% GUION. Fuentes: doc/SoTA/comparative_analysis.md y doc/memoria/materiales.md seccion 6.
% REGLA: ninguna cita sin DOI verificado. Los preprints se citan indicando que lo son.
% Excluidos por no citables: Wang/Zhengzhou, Stirbu, Al Farib.


\section{Mitigación clásica de errores}
\label{sec:qem_clasico}

\subsection{Zero Noise Extrapolation (ZNE)}
Para corregir los errores, ZNE propone amplificar el ruido de forma controlada y luego extrapolar hacia atrás, imaginando cómo sería el resultado en un escenario de ruido cero[cite: 3]. El gran problema es \textbf{su coste: nos obliga a lanzar varias ejecuciones extra por cada resultado que queremos obtener}[cite: 1, 3].

\subsection{Probabilistic Error Cancellation (PEC)}
PEC toma otro camino y se basa en el muestreo de una cuasi-distribución[cite: 3]. Es un método insesgado, pero tiene una trampa terrible en la práctica: su coste computacional \textbf{crece exponencialmente} a medida que aumenta el ruido del sistema[cite: 3].

\subsection{Mitigación de lectura (M3)}
M3 se centra específicamente en los errores de lectura, y lo hace invirtiendo la matriz de asignación sin llegar a construirla explícitamente de forma completa[cite: 1].

\subsection{El denominador común}
¿Qué tienen en común todos estos métodos clásicos? Que todos gastan ejecuciones adicionales sobre un recurso que ya es bastante escaso de por sí, y, sobre todo, que todos actúan \textbf{después} de haber ejecutado el circuito[cite: 1].

\section{Aprendizaje automático aplicado a la mitigación}
\label{sec:ml_qem}

\subsection{Modelos sustitutos entrenados sobre resultados ruidosos}
El estudio de Liao et al. (2024) es nuestra referencia central aquí[cite: 2, 3]. Demostraron que podemos entrenar redes que aprendan a corregir directamente a partir del valor ruidoso, ahorrándonos el enorme coste extra de ejecuciones que exigen métodos como ZNE[cite: 3]. \textbf{El dato que este trabajo aprovecha:} en sus pruebas de ablación demostraron que \textbf{Random Forest} es el mejor modelo tabular de referencia[cite: 2, 3]. Precisamente por eso está justificado incluirlo como uno de los pilares en la comparativa del capítulo~\ref{cap:comparativa}[cite: 1, 4].

\subsection{Arquitecturas de grafo}
En el caso de GTraQEM (Bao et al., ICLR 2025), plantearon un \textit{Graph Transformer} \textbf{sin paso de mensajes} pero apoyado en un \textbf{nodo virtual}[cite: 2, 4]. De aquí saco exactamente la idea del nodo virtual que sí uso en este trabajo, pero \textbf{su arquitectura no se reproduce}[cite: 2, 4]: el porqué descarté su diseño lo desarrollo a fondo en el capítulo~\ref{cap:comparativa}[cite: 4].

Por otro lado, Liu et al. (2026) montaron una red neuronal de grafos (GNN) sobre el grafo acíclico dirigido (DAG) usando un vector por nodo muy parecido al nuestro (ellos usan 31 dimensiones y yo 25)[cite: 2]. Esto nos viene genial para justificar por qué apostamos por GNN en lugar de por redes convolucionales (CNN)[cite: 2].

\subsection{El techo actual de rendimiento}
Actualmente, el rey en rendimiento es QEMFormer (Bao et al., ICML 2025) apoyado en el banco de pruebas QEM-Bench[cite: 2]. \textbf{Y aquí está el contraste central de esta memoria:} QEMFormer necesita que le des el valor ruidoso \textbf{como entrada del modelo}[cite: 2, 4]. Es decir, tienes que ejecutar el circuito en la máquina antes de poder mitigarlo[cite: 1, 2]. Además, como codifican los qubits con un formato \textit{multi-hot}, el modelo se queda atado obligatoriamente al tamaño exacto de qubits con el que se entrenó[cite: 4].

\section{El paradigma pre-ejecución}
\label{sec:preejecucion}

\subsection{Heurísticas: mapomatic}
Nation \& Treinish (2023) presentan la herramienta oficial de IBM, mapomatic[cite: 2]. Lo que hace es puntuar \textbf{layouts alternativos del mismo circuito} basándose en la calibración del día para quedarse con el mejor camino, pero sin llegar a usar aprendizaje automático[cite: 2].

\textbf{Esta distinción hay que hacerla muy explícita:} mapomatic \textbf{elige} el mejor mapa posible; este trabajo, en cambio, \textbf{estima una magnitud} concreta para poder aplicar una corrección matemática[cite: 2]. Son dos objetivos totalmente distintos[cite: 2].

\subsection{Ranking aprendido: Q-CTRL}
Hartnett et al. (2024) usan \textit{machine learning} para ordenar circuitos equivalentes antes de mandarlos a ejecutar[cite: 2]. Sin embargo, su modelo solo predice el \textbf{orden} (cuál irá mejor o peor), pero no predice la magnitud matemática, así que tampoco nos sirve para corregir el resultado final[cite: 2].

\section{Variabilidad temporal del hardware}
\label{sec:drift}
Hirasaki et al. (2023) demostraron algo vital: las tasas de error de los qubits no empeoran poco a poco, sino que cambian de forma \textbf{escalonada} y brusca[cite: 2, 4]. Esta es la evidencia empírica que justifica por qué este dataset tiene un eje temporal, y es lo que invalidó por completo la idea que tenía en el diseño original de simular una deriva lineal del ruido[cite: 4].

Además, Huo et al. (2026) confirman esta misma variabilidad tremenda, pero vista desde los resultados operacionales que sufre el usuario final[cite: 2].

\section{Paradigmas alternativos}
\label{sec:alternativas}
Existen otros enfoques, como Q-Cluster (Patil et al., 2025), que aborda una mitigación \textbf{no supervisada}, es decir, sin usar etiquetas ideales[cite: 2]. En la misma línea, Liao et al. (2025) proponen mitigación sin \textit{ground truth}[cite: 2]. Ambos trabajos los menciono también en el apartado de limitaciones, porque nuestro pipeline \textbf{sí} es totalmente dependiente de tener un valor exacto que podamos simular[cite: 2].

\section{El hueco}
\label{sec:gap}
En definitiva, nadie está prediciendo la \textbf{magnitud} de la degradación \textbf{antes de ejecutar} combinando la estructura del circuito y la telemetría del chip[cite: 1, 2]. Nuestro modelo viene a cubrir esto con cuatro diferenciadores concretos:
\begin{enumerate}
    \item Opera puramente en pre-ejecución: la entrada del modelo nunca incluye el resultado medido[cite: 2].
    \item No codifica rígidamente la identidad del qubit, así que generaliza de forma natural a circuitos de cualquier tamaño[cite: 2, 4].
    \item Modela el \textit{drift} temporal usando calibraciones reales de 42 días[cite: 2, 4].
    \item Evalúa circuitos algorítmicos complejos como QFT de forma sistemática, algo que ningún otro trabajo de ML-QEM hace[cite: 2].
\end{enumerate}